\subsection{50. Novel Bile Acid Biosynthetic Pathways Enriched in the Centenarian Microbiome}
\textbf{PMID:} \texttt{34325466}\quad\textbf{Authors:} Sato et al.\ (2021)\quad\textbf{Journal:} Nature 599, 458--464\quad\textbf{Domain:} Metagenomics / Bile Acid Metabolism\\
\scorebadge{Coverage}{4}\quad\scorebadge{Agreement}{6}\quad\textbf{Total:} 10/20\par\smallskip
\noindent\textbf{Coverage rationale.} Spot-check: 10 centenarian + 10 elderly samples at 3--7\% read depth (500K--1M of 13--20M total reads). Reference genomes from 4 Odoribacteraceae isolates + \textit{P.\ merdae} used for read mapping. 21 5$\alpha$-reductase and 10 3$\beta$-HSDH CDS queried. No MAG assembly, no metabolomics.\par
\noindent\textbf{Agreement rationale.} 5$\alpha$-reductase genes show 1.30$\times$ enrichment in centenarians (directionally consistent). \textit{A.\ finegoldii} 5AR enriched 5.29$\times$; \textit{P.\ merdae} 4.00$\times$. Overall Odoribacteraceae abundance not significantly different ($p=0.71$). At 3--7\% read depth, organism-specific genes are too sparse for statistical power.\par\smallskip
\noindent\textbf{Verdict: SPOT-CHECK.} Directionally consistent but far too shallow for definitive replication. Full-depth assembly pipeline required.

